\begin{abstract}
Aerial--ground cooperative perception improves scene understanding through information exchange across agents, yet intermediate-fusion detection outputs can still omit useful candidates and rank their quality inadequately. We propose GeoRank, which uses geometry-disagreement-aware ranking to strengthen late supplementation of cooperative outputs with independent vehicle predictions. GeoRank associates two sets of candidates at the receiver and incorporates center disagreement, matching cost, confidence relationships, and source status into final scoring. It supports both explicit disagreement penalties and lightweight learned rescoring, without retraining the underlying detection networks for ranking. The independent vehicle branch uses local observations, while supplementation and ranking reuse the same remote inputs, introducing no additional remote message payload relative to the cooperative baseline. On the complete Griffin-55m validation set, the explicit multiplicative variant achieves \textbf{39.87 vehicle AP and 44.68 vehicle AMOTA}, numerically exceeding all intermediate-fusion methods reported in Table 3 of the Griffin benchmark and surpassing the strongest reported method, CoopTrack, by \textbf{3.47 and 4.48 percentage points}, respectively. With identical branch inputs, linear and multiplicative disagreement scoring improve the base fusion AP from 39.26 to \textbf{39.86 and 39.87}. With candidates and box geometry further fixed, adding center disagreement alone increases AP on a five-scene holdout within the validation set from 38.70 to \textbf{39.49 (+0.79)}. Learned GeoRank reaches \textbf{40.12 (+1.42)} and improves AP by \textbf{2.38 points} at the strict 0.5 m matching tolerance. Module ablations, eight feature configurations, and scene-level analyses support the ranking value of geometric disagreement. Explicit penalties slightly reduce AMOTA relative to base fusion, indicating that detection ranking and tracking operating points still require coordination. These results show that receiver-side candidate relationship modeling can improve late supplementation in 55 m aerial--ground cooperative perception while preserving the upstream message payload.

\textbf{Keywords:} aerial--ground cooperative perception; 3D object detection; late fusion; geometric disagreement; candidate rescoring.
\end{abstract}
